\chapter{Derived affine intersections and cotangent complexes}
\label{ch:derived-affines}

Let \(k\) be a field of characteristic zero. In this chapter algebras
are commutative dg \(k\)-algebras in cochain degrees \(\leq0\),
with \(ab=(-1)^{|a||b|}ba\). An ordinary \(k\)-algebra is placed
in degree zero. All constructions use graded commutativity and
specified dg algebra maps. The affine convention reverses algebra
maps. We call the object represented by \(B\) a derived affine and
regard quasi-isomorphic algebras as equivalent. We prove the affine
calculations underlying this notation, following the scope of
\cite[\S\S2.2 and 4.1]{toen}.

\section{Semifree resolutions of commutative dg algebras}

\begin{definition}[Semifree algebra extension]\label{res:A01}
A semifree extension of \(A\) is a graded-commutative polynomial
extension on homogeneous variables adjoined in stages, with
the differential of a variable in the algebra at earlier stages.
The free graded-commutative algebra is polynomial in even variables
and exterior in odd variables. If prescribed differentials are
cycles of the required degrees, the unique graded derivation extending
them is a differential.
\end{definition}
\begin{proof}
The Koszul sign rule gives a well-defined derivation on the free
graded-commutative algebra. In particular for an odd variable \(x\),
\(d(x^2)=d(x)x-xd(x)=0\) since \(d(x)\) is even.
The cross terms in applying a degree-one derivation twice cancel.
Thus \(d^2\) vanishes on every monomial if it vanishes on the base
and each variable. Each specified cycle ensures that condition.
\end{proof}

\begin{theorem}[Semifree resolution and lifting]\label{res:A02}
Every map \(A\to B\) admits a factorization
\(A\to P\twoheadrightarrow B\) with \(P\) semifree over \(A\)
and the second arrow a degreewise surjective quasi-isomorphism.
A map from a semifree extension lifts through any such surjective
quasi-isomorphism of \(A\)-algebras.
\end{theorem}
\begin{proof}
In degree zero adjoin cycles mapping to all elements of \(B^0\).
For each \(b\in B^n\), \(n<0\), first adjoin a cycle \(v_b\)
in degree \(n+1\) mapping to \(db\), then \(u_b\) in degree \(n\)
with \(du_b=v_b\), mapping to \(b\). This makes the map onto \(B\)
surjective in every degree. At the next stage, for every homogeneous
cycle \(z\) in the kernel adjoin a generator of degree \(|z|-1\)
with differential \(z\) and image zero. Repeat countably many times.
Every element is a finite polynomial using finitely many stages.
Thus every kernel cycle in the union occurred at a stage and became
a boundary at the next. The kernel is acyclic and the surjection is
a quasi-isomorphism.

For lifting through \(C\twoheadrightarrow D\), choose the lift of a
generator \(x\) after all generators in \(dx\) have been lifted.
A candidate lift \(c\) has \(dc-\widetilde{dx}\) a cycle in the kernel.
Subtract a primitive of that cycle to get the required lift. Proceed
through the stages. In degree zero the candidate is already a cycle
because the target is nonpositive. The construction respects the base
and is defined on the union.
\end{proof}

\begin{lemma}[K-flatness of a semifree algebra]\label{res:A03}
A semifree \(P\) over \(A\), regarded as a dg \(A\)-module,
is K-flat.
\end{lemma}
\begin{proof}
At one adjoining stage use the increasing filtration by total degree
in the new variables. Differentiating such a variable removes it and
gives a coefficient from the previous stage. Each graded quotient of
the filtration is therefore a direct sum of shifts of the previous
stage as a dg module. The inclusions are split as graded modules by
the monomial bases. Tensor with an acyclic dg \(A\)-module. Starting
from \(A\), the quotients after tensoring are acyclic and the graded
splittings preserve exactness of successive sequences after tensoring.
The long exact sequence and then an exact filtered union prove
acyclicity at that adjoining stage. Induct through the stages, using
filtered exactness at their unions.
\end{proof}

\section{Derived base change and intersection}

\begin{proposition}[Tensor model for a derived affine base change]\label{res:A04}
For \(A\to B\) and \(A\to C\), a model of the derived tensor product
is \(P\otimes_A C\) for a semifree surjection \(P\to B\) as in
Result~\ref{res:A02}. A change of models gives a zigzag of algebra
quasi-isomorphisms. Its zeroth cohomology is
\[
 H^0(P\otimes_A C)
   \cong H^0(B)\otimes_{H^0(A)} H^0(C).
\]
\end{proposition}
\begin{proof}
Choose semifree surjective resolutions of both \(B\) and \(C\).
Tensoring the two resolutions computes the derived tensor product of
dg \(A\)-modules by Result~\ref{res:A03}. Replacing the second one
by \(C\) is then a quasi-isomorphism by K-flatness of \(P\).
For two choices for \(B\), lift the augmentation from one to the
other by Result~\ref{res:A02}; the lifted map is a quasi-isomorphism
by the comparison to \(B\). Tensor with the semifree model for \(C\)
and use its K-flatness. This gives the stated zigzag of maps of
commutative dg algebras, and the argument is symmetric in the two
inputs. We use this construction as the derived affine fibre product
model.

Since all degrees are nonpositive, degree zero of the tensor is
generated by \(P^0\otimes_{A^0}C^0\), and its boundaries are exactly
the images obtained by using \(dP^{-1}\) in the first factor or
\(dC^{-1}\) in the second. The balancing relations in degree zero
come from \(A^0\). Its boundaries act as zero after the two
quotients, so the result factors over \(H^0(A)\), giving the formula.
\end{proof}

\begin{example}[The Koszul model of a derived zero locus]\label{res:A05}
Let \(A\) be ordinary and let \(f_1,\ldots,f_r\in A\).
Regard \(A\) as an algebra over \(k[t_1,\ldots,t_r]\) by \(t_i\mapsto f_i\).
The derived fibre over the zero section has algebra
\[
     A[e_1,\ldots,e_r],\qquad |e_i|=-1,\qquad de_i=f_i .
\]
If \(f_i\) form an \(A\)-regular sequence then this algebra resolves
\(A/(f_1,\ldots,f_r)\). In that case the derived self-intersection
of the closed locus over \(A\) has exterior generators in degree
\(-1\) and zero differential over \(A/(f_1,\ldots,f_r)\).
\end{example}
\begin{proof}
The analogous algebra with differential \(de_i=t_i\) is the Koszul
resolution of \(k\) over the polynomial ring, because the variables
are regular. The exactness is the Koszul resolution calculation
\cite[RK-07]{ca-core}. It is semifree, so tensoring with \(A\)
gives the derived fibre by Result~\ref{res:A04}.
For a regular sequence in \(A\), the same Koszul exactness applies
over \(A\). The displayed algebra is then a semifree resolution of
the quotient over \(A\). Tensor with the quotient, which sends
each \(f_i\) to zero. For example the derived self-intersection of
the origin in \(\mathbb A^1_k\) has \(H^0=k\) and \(H^{-1}=k e_1\).
\end{proof}

\section{The cotangent calculation}

Use \(\partial\) for a K\"ahler differential and \(d\) for the
cochain differential. For a dg \(P\)-module, the usual graded Hom
complex is formed using the Koszul rule for the degree of a
homogeneous linear map.

\begin{lemma}[Differentials on a semifree algebra]\label{res:A06}
For \(A\to P\) semifree, the dg module \(\Omega_{P/A}\) is free
as a graded \(P\)-module on \(\partial x\) for the variables \(x\),
where \(|\partial x|=|x|\) and \(d(\partial x)=\partial(dx)\).
It is a semifree \(P\)-module. Its base change to any dg \(P\)-algebra
is again semifree. The complex
\(\operatorname{Hom}^{\bullet}_P(\Omega_{P/A},M)\) identifies with
graded \(A\)-derivations \(P\to M\) with differential the graded
commutator with \(d\).
\end{lemma}
\begin{proof}
The graded Leibniz relation extends arbitrary values on the free
generators uniquely to a graded derivation, and it respects
the symmetric and exterior relations. Hence the graded module of
differentials has the stated basis. The relation requiring
\(\partial\) to commute with the cochain differential determines
the formula for \(d\) on that basis. The derivative of \(dx\) uses
only earlier generators, so the corresponding filtration gives
semifreeness and is preserved by base change. The Hom--derivation
identification now commutes with its differential by evaluating
the commutators on the same basis.
\end{proof}

\begin{lemma}[A relative polynomial path comparison]\label{res:A07}
Let \(Q\twoheadrightarrow B\) be a surjective quasi-isomorphism of
nonpositive dg algebras. Let \(A\to P\) be semifree with a map
\(P\to B\).
Any two \(A\)-algebra maps \(P\to Q\) over \(B\) are joined by a
polynomial dg path over \(B\). Consequently their induced maps
\[
 B\otimes_P\Omega_{P/A}\longrightarrow B\otimes_Q\Omega_{Q/A}
\]
are chain homotopic as maps of dg \(B\)-modules.
\end{lemma}
\begin{proof}
Let \(k[t,dt]\) have \(|t|=0\), \(|dt|=1\) and differential
\(d(t)=dt\), with \(dt\) exterior. Endpoints set \(t=0,1\) and
\(dt=0\). The inclusion of constants and either evaluation are
mutual homotopy inverses as complexes by polynomial integration.
On a homogeneous coefficient \(y\), the integration of
\(y\,g(t)dt\) has sign \((-1)^{|y|}\).
The identity \(dH+Hd=\operatorname{ev}_1-\operatorname{ev}_0\)
for this integration follows by the fundamental formula for a
polynomial derivative and the opposite signs on the two terms
containing \(d y\). Integers are invertible in \(k\).

Take the relative path algebra
\[
 T=\tau_{\leq0}\bigl(Q[t,dt]\times_{B[t,dt]} B\bigr),
\]
where the truncation keeps the kernel of \(d\) in degree zero and
all degrees below zero. This is a subalgebra. Write \(K\) for the
acyclic kernel of \(Q\to B\). Before truncation, evaluation at zero
has section given by constants and kernel
\(K\otimes_k\ker(k[t,dt]\to k)\) for evaluation at zero.
This kernel is contractible using integration from zero to \(t\).
The good truncation preserves the cohomology comparison to \(Q\).
The projection \(Q\times_B Q\to Q\) is also a quasi-isomorphism
since its kernel is \(K\). Therefore
\(\operatorname{ev}_{01}:T\to Q\times_B Q\) is a quasi-isomorphism.
It is surjective in negative degrees using linear interpolation
of endpoints. In degree zero, if endpoints differ by \(w\in K^0\),
pick \(z\in K^{-1}\) with \(dz=w\). Then the path
\(q_0+tw-z\,dt\) is closed and has the required endpoints.
Use the lifting part of Result~\ref{res:A02} for \(P\) to get the
path.

To compare on differentials, differentiate in the \(Q\)-variable
of the path, leaving \(t,dt\) constant. This maps differentials
into \(\Omega_{Q/A}[t,dt]\). After extension along \(Q\to B\), the
original path has constant image in \(B[t,dt]\). The derivative
therefore induces a dg \(B\)-linear map into
\((B\otimes_Q\Omega_{Q/A})[t,dt]\). Its endpoint maps are the two
maps in the statement. Integration gives a graded \(B\)-linear
chain homotopy between them.
\end{proof}

\begin{theorem}[The affine cotangent complex]\label{res:A08}
For a semifree surjective resolution \(P\to B\) over \(A\), the
object
\[
       L_{B/A}=B\otimes_P\Omega_{P/A}
\]
is independent of the resolution up to canonical isomorphism in
the homotopy category of dg \(B\)-modules. It is represented by
a K-projective, K-flat dg \(B\)-module.
\end{theorem}
\begin{proof}
For two resolutions lift the maps over \(B\) in both directions
by Result~\ref{res:A02}. Their composites and the identities
agree after augmentation. Result~\ref{res:A07} proves that their
differential maps are mutual homotopy inverses and independent
of the chosen lifts. Composing three comparisons respects these
homotopy classes, giving the canonical isomorphisms.
The dg module in the statement is semifree by
Result~\ref{res:A06}. A semifree dg module is K-flat by its
filtration with shifted free quotients. To see K-projectivity,
a cocycle map from it into an acyclic dg module can be made
nullhomotopic successively on generators. On a new generator,
subtract the already specified homotopy on its differential.
The resulting prescribed boundary equation has a cycle on
its left, so acyclicity supplies a primitive. The Koszul sign
determines the module-linear extension of the homotopy.
Apply this in every degree by shifting the target. Thus the
Hom complex into an acyclic dg module is acyclic.
\end{proof}

\begin{theorem}[Transitivity and base change of the cotangent complex]\label{res:A09}
For \(A\to B\to C\) as above there is a natural triangle in \(D(C)\)
\[
 C\otimes_B^{\mathbf L}L_{B/A}\longrightarrow L_{C/A}
 \longrightarrow L_{C/B}\longrightarrow
 (C\otimes_B^{\mathbf L}L_{B/A})[1].
\]
For \(A\to A'\), if \(P\to B\) is semifree and
\(B'=A'\otimes_A P\) is the derived base-change model, then
\[
       L_{B'/A'}\simeq B'\otimes_P \Omega_{P/A}.
\]
In the flat ordinary base-change case, the second formula identifies
with extension of scalars on \(L_{B/A}\) along \(B\to A'\otimes_A B\).
\end{theorem}
\begin{proof}
Choose \(P\to B\) semifree over \(A\), then resolve \(C\) by \(Q\)
semifree over \(P\) along \(P\to B\to C\), using
Result~\ref{res:A02}. Concatenation makes \(Q\) semifree over \(A\).
The variable bases give a graded-split exact sequence of dg modules
\[
 0\longrightarrow Q\otimes_P\Omega_{P/A}
 \longrightarrow \Omega_{Q/A}
 \longrightarrow \Omega_{Q/P}\longrightarrow0 .
\]
The first submodule is a subcomplex since the differential of
a \(P\)-variable uses only \(P\)-variables. The quotient is specified
by setting \(\partial x=0\) for each \(P\)-variable \(x\).
Tensor with \(C\) over \(Q\); graded splitness preserves exactness.
The first resulting term equals
\(C\otimes_B(B\otimes_P\Omega_{P/A})\), with the factor in parentheses
K-flat by Result~\ref{res:A08}. The middle term computes \(L_{C/A}\).
For the last term \(B\otimes_P Q\) is semifree over \(B\) and
quasi-isomorphic to \(C\): the map from \(Q\) is a quasi-isomorphism
by Result~\ref{res:A03}, and \(Q\to C\) is a quasi-isomorphism.
The map \(B\otimes_PQ\to C\) is surjective since \(Q\to C\) is.
Relative differentials commute with extension of scalars by the
variable formula. Thus the last term computes \(L_{C/B}\).
The exact sequence supplies the triangle. To compare two resolved
diagrams, lift first on \(P\) and then on the variables of \(Q\),
over the prescribed maps on \(B,C\). Such a comparison commutes
with the injection and quotient maps on differentials. For two
choices of comparison to \(P'\to Q'\), choose a relative path
between the two maps on \(P\) over \(B'\) by
Result~\ref{res:A07}. The map \(P'\to Q'\) carries it to the
relative path algebra of \(Q'\to C'\), which now has a \(P\)-algebra
structure. The endpoint map to \(Q'\times_{C'}Q'\) is a surjective
quasi-isomorphism. Lift through it using the semifree \(Q\)-cells
over \(P\). This yields a path of diagram comparisons with the
chosen path on \(P\). Differentiation and integration as in
Result~\ref{res:A07} give compatible homotopies on the short exact
sequences, proving independence and naturality in the derived
category (with extension of scalars on the target modules).

Finally \(A'\otimes_A P\) is itself semifree over \(A'\);
the relative variables give the claimed expression on the same
generators and differential. For flat ordinary base change the
map \(A'\otimes_A P\to A'\otimes_A B\) is a surjective
quasi-isomorphism. Using this resolution and the associativity
of tensor gives the final identification.
\end{proof}

\begin{proposition}[Square-zero tangent calculation]\label{res:A10}
Let \(M\) be a nonpositive dg \(B\)-module. Augmented maps
\(P\to B\ltimes M\) of \(A\)-algebras over the fixed augmentation
\(P\to B\) are naturally the degree-zero cycles in
\(\operatorname{Hom}_B^\bullet(L_{B/A},M)\) for the model of
Result~\ref{res:A08}. Their classes modulo polynomial dg homotopy
over \(B\) are \(H^0\mathrm{RHom}_B(L_{B/A},M)\).
All groups \(H^i\mathrm{RHom}_B(L_{B/A},M)\) are computed by the
same Hom complex.
\end{proposition}
\begin{proof}
The multiplication on \(B\ltimes M\) has square-zero ideal \(M\).
The difference of an augmented map from the map with zero
\(M\)-component is therefore an \(A\)-derivation by the Leibniz rule.
Commutation with the differential is its degree-zero cocycle
condition. Result~\ref{res:A06} gives the correspondence, and
K-projectivity in Result~\ref{res:A08} computes the derived Hom.

For the assertion on classes, a polynomial path is a closed
degree-zero derivation with values in \(M[t,dt]\) over the
constant augmentation, with a good truncation if necessary.
Integration as in Result~\ref{res:A07} makes the endpoints differ
by a Hom boundary. Conversely, if two cocycles \(f_0,f_1\) have
\(f_1-f_0=d_{\rm Hom}h\) for \(h\) of degree \(-1\), the path
\(f_0+t(f_1-f_0)-h\,dt\), using the tensor sign convention, is
closed of degree zero. Here the expression means the map into
\(M\otimes_k k[t,dt]\) obtained by tensoring graded maps with the
specified forms. It is well-defined by the graded tensor sign
on homogeneous arguments. Its degree-zero values on closed
elements are closed, so it factors through the good truncation.
It therefore gives the required dg algebra path. This proves
the quotient identification.
\end{proof}

\section{Complete intersections and compatibility with descent}

\begin{example}[The conormal complex and its tangent space]\label{res:A11}
Let \(B=S/(f_1,\ldots,f_r)\) with \(S=k[x_1,\ldots,x_s]\) and
\(f_i\) regular in \(S\). Then \(L_{B/k}\) is represented by
\[
    [\,B^r\xrightarrow{(\partial f_i/\partial x_j)}B^s\,]
\]
in degrees \(-1,0\), with column \(i\) the derivative of \(f_i\).
In particular \(H^1\mathrm{RHom}_B(L_{B/k},B)\) is the first-order
deformation space of the marked algebra in
\cite[T04]{ca-topics}.
\end{example}
\begin{proof}
Use the semifree Koszul algebra over \(S\) of Result~\ref{res:A05}
and regard the variables \(x_j\) as its first cells over \(k\).
The differential on \(\partial e_i\) is \(\partial f_i\);
the differentials on \(\partial x_j\) vanish. Tensor with \(B\)
to get precisely the two-term complex.
Taking Hom into \(B\) gives in degree one the cokernel of
the transposed Jacobian. The explicit square-zero presentation
calculation in \cite[T02--T04]{ca-topics} identifies that cokernel
with the deformation classes over the dual numbers.
\end{proof}

\begin{proposition}[Faithfully flat locality of the affine tangent calculation]\label{res:A12}
Let \(A\) be an ordinary \(k\)-algebra, \(A\to C\) faithfully
flat, and \(B\) an ordinary \(A\)-algebra. Assume the affine
cotangent complex \(L_{B/A}\) is represented by a bounded complex
\(P_L\) of finite projective \(B\)-modules. Set
\(C_p=C^{\otimes_A(p+1)}\) and \(B_p=B\otimes_A C_p\).
For any complex \(M\) of \(B\)-modules there is a comparison
\[
 \mathrm{RHom}_B(L_{B/A},M)
 \simeq \operatorname{Tot}^{\Pi}_{p\geq0}
     \mathrm{RHom}_{B_p}(L_{B_p/C_p}, B_p\otimes_B M).
\]
When all local tangent calculations are made in the ordinary
degree-zero case, this compares their classes and higher
cohomology by the indicated full total complex.
\end{proposition}
\begin{proof}
Base change \(B\to B\otimes_A C\) is faithfully flat.
Each \(C_p\) is flat over \(A\). The base-change part of
Result~\ref{res:A09} uses a single semifree model over \(A\)
and its compatible scalar extensions to the whole Amitsur diagram.
Thus \(B_p\otimes_B P_L\) represents \(L_{B_p/C_p}\).
A comparison from \(P_L\) to the cofibrant cotangent model can
be chosen as a chain map since \(P_L\) is K-projective, and its
flat scalar extensions commute with the cofaces. Apply
Result~\ref{res:F08} over \(B\) to \(P_L\) and \(M\), interpreting
the Hom complexes using these common models.
\end{proof}

The affine calculations here use a characteristic-zero commutative
dg model. In general, derived stacks involve functors of points,
homotopy-coherent descent and representability. These broader
constructions, as well as the all-characteristic simplicial
setting, are developed in \cite[\S\S1.2.12 and 2.2.3]{hag} and
\cite[\S\S2 and 4]{toen}.
